% Biên dịch độc lập: tectonic docs/theory/part-b.tex (hoặc XeLaTeX hai lượt).
\documentclass[11pt,a4paper]{article}
\usepackage[margin=22mm,headheight=14pt]{geometry}
\usepackage{fontspec}
\setmainfont{texgyretermes}[Extension=.otf,UprightFont=*-regular,BoldFont=*-bold,ItalicFont=*-italic,BoldItalicFont=*-bolditalic]
\usepackage{amsmath,amssymb}
\usepackage{fancyhdr}
\usepackage[unicode,hidelinks]{hyperref}
\usepackage{bookmark}

\hypersetup{pdftitle={Phần B — Chứng minh lý thuyết},pdflang={vi-VN}}
\setlength{\parindent}{0pt}
\setlength{\parskip}{5pt plus 1pt minus 1pt}
\setlength{\emergencystretch}{2em}
\linespread{1.08}
\raggedbottom
\renewcommand{\thesection}{B.\arabic{section}}
\renewcommand{\contentsname}{Mục lục}
\renewcommand{\refname}{Tài liệu tham khảo}
\renewcommand{\headrulewidth}{0.3pt}
\pagestyle{fancy}
\fancyhf{}
\fancyhead[L]{\small Đề tài 9 — Chứng minh lý thuyết}
\fancyhead[R]{\small Bản nháp}
\fancyfoot[C]{\thepage}

\DeclareMathOperator{\Var}{Var}
\DeclareMathOperator{\Cov}{Cov}
\DeclareMathOperator{\rank}{rank}
\DeclareMathOperator{\im}{im}
\DeclareMathOperator{\diag}{diag}
\DeclareMathOperator{\softmax}{softmax}
\DeclareMathOperator{\KL}{KL}
\DeclareMathOperator{\supp}{supp}
\DeclareMathOperator*{\argmin}{arg\,min}
\newcommand{\E}{\mathbb{E}}
\newcommand{\R}{\mathbb{R}}
\newcommand{\PE}{\mathrm{PE}}
\newcommand{\Attn}{\operatorname{Attn}}
\newcommand{\PPL}{\mathrm{PPL}}
\newcommand{\pref}{\pi_{\mathrm{ref}}}
\newcommand{\pdata}{p_{\mathrm{data}}}

\begin{document}
{\LARGE\bfseries Phần B — Chứng minh lý thuyết\par}
{\large Kiến trúc Transformer và tinh chỉnh mô hình ngôn ngữ\par}
\textbf{Nhóm:} Cao Văn Hà (trưởng nhóm), Hà Mạnh Thế, Trần Duy Hưng.\par
\textbf{Phạm vi:} Tài liệu chứng minh sáu mệnh đề B.1--B.6 của Đề tài 9 (đề gốc ngày 12/11/2026, tr.~18) \cite{de}. Các chứng minh không cần đầu vào của Phần C. Những ví dụ số trong tài liệu chỉ dùng để minh họa. Mọi logarit là logarit tự nhiên.\par
\textbf{Trạng thái:} Bản nháp có hỗ trợ AI, chưa được nhóm xác nhận là bản nộp.
\tableofcontents

\section{Phương sai của tích vô hướng và hệ số chia tỷ lệ}
\label{sec:scale}
\textbf{Mệnh đề.} Cho $d_k\ge 1$ và $q,k\in\R^{d_k}$. Giả sử toàn bộ $2d_k$ thành phần của $q$ và $k$ độc lập lẫn nhau, mỗi thành phần có kỳ vọng $0$ và phương sai $1$. Khi đó $\Var(q^\top k)=d_k$. Hệ số $1/\sqrt{d_k}$ giữ phương sai của mỗi điểm chú ý đã chia tỷ lệ bằng $1$ trong mô hình giả thiết này \cite[§3.2.1, chú thích 4]{attention}.

\subsection*{Tính phương sai}
Đặt $Z_i=q_i k_i$, với $1\le i\le d_k$. Từ $\Var(q_i)=\E[q_i^2]-(\E[q_i])^2$, suy ra $\E[q_i^2]=1$; tương tự $\E[k_i^2]=1$. Vì $q_i$ và $k_i$ độc lập,
\[
  \E[Z_i]=\E[q_i]\E[k_i]=0,\qquad
  \E[Z_i^2]=\E[q_i^2]\E[k_i^2]=1,
\]
nên $\Var(Z_i)=1$. Với $i\ne j$, $Z_i$ và $Z_j$ là các hàm của hai nhóm biến độc lập $(q_i,k_i)$ và $(q_j,k_j)$, do đó $\Cov(Z_i,Z_j)=0$. Vì $q^\top k=\sum_{i=1}^{d_k} Z_i$,
\begin{align*}
  \E[q^\top k]&=0,\\
  \Var(q^\top k)
    &=\sum_{i=1}^{d_k}\Var(Z_i)+2\sum_{1\le i<j\le d_k}\Cov(Z_i,Z_j)=d_k,\\
  \Var\!\left(\frac{q^\top k}{\sqrt{d_k}}\right)&=\frac{1}{d_k}\Var(q^\top k)=1.
\end{align*}

\subsection*{Vì sao điều này giúp giảm bão hòa softmax?}
Xét $m\ge2$ điểm chú ý $z_1,\ldots,z_m$ trong một hàng. Viết $s_i=e^{z_i}/\sum_{j=1}^m e^{z_j}$. Dùng quy tắc đạo hàm thương, với $\delta_{ij}=1$ nếu $i=j$ và $0$ nếu $i\ne j$, ta có
\[
  \frac{\partial s_i}{\partial z_j}
  =\frac{\delta_{ij}e^{z_i}}{\sum_a e^{z_a}}
   -\frac{e^{z_i}e^{z_j}}{(\sum_a e^{z_a})^2}
  =s_i(\delta_{ij}-s_j).
\]
Nếu điểm lớn nhất $z_a$ hơn mọi điểm khác ít nhất $g>0$ thì
\[
  s_a=\frac{1}{1+\sum_{j\ne a} e^{-(z_a-z_j)}}
    \ge \frac{1}{1+(m-1)e^{-g}}.
\]
Giữ $m$ cố định và cho $g$ tăng, $s_a$ tiến đến $1$, các $s_j$ còn lại tiến đến $0$ và toàn bộ Jacobian của softmax tiến đến $0$. Khi đó trọng số chú ý ít thay đổi trước nhiễu nhỏ ở các điểm số. Bão hòa phụ thuộc vào \emph{chênh lệch} điểm, không phải việc cộng một hằng số chung vào cả hàng.

Với các tích vô hướng thô, độ lệch chuẩn của từng điểm là $\sqrt{d_k}$ dưới giả thiết trên. Chia cả hàng cho $\sqrt{d_k}$ đưa độ lệch chuẩn của mỗi điểm về $1$ và thu nhỏ các chênh lệch theo cùng hệ số nếu $d_k>1$. Kết quả phương sai giải thích lý do chọn thang điểm; nó không khẳng định mọi hàng đều không bão hòa, hay các tọa độ của query và key thực tế luôn độc lập.

\section{Tính đẳng biến hoán vị của self-attention}
\label{sec:perm}
\textbf{Mệnh đề.} Xét self-attention không mã hóa vị trí, không mặt nạ và dùng cùng trọng số chiếu cho mọi hàng. Với $X\in\R^{n\times d_{\mathrm{model}}}$, mỗi hàng là biểu diễn một token. Nếu $P\in\R^{n\times n}$ là ma trận hoán vị thì $\Attn(PX)=P\Attn(X)$. Đây là \emph{đẳng biến}: đầu ra đổi thứ tự theo đầu vào, không phải đầu ra bất biến.

\subsection*{Theo dõi phép hoán vị đến ma trận điểm}
Dùng $W_Q,W_K\in\R^{d_{\mathrm{model}}\times d_k}$ và $W_V\in\R^{d_{\mathrm{model}}\times d_v}$, với $n,d_{\mathrm{model}},d_k,d_v>0$. Đặt
\[
  Q=XW_Q,\quad K=XW_K,\quad V=XW_V,\qquad
  S=\frac{QK^\top}{\sqrt{d_k}}\in\R^{n\times n}.
\]
Hai ma trận $Q,K$ có kích thước $n\times d_k$; $V$ có kích thước $n\times d_v$. Theo định nghĩa chú ý tích vô hướng có chia tỷ lệ \cite[§3.2.1]{attention}, $\Attn(X)=\softmax_{\mathrm{row}}(S)V$, trong đó softmax áp dụng riêng cho từng hàng. Với đầu vào $PX$, các phép chiếu cho $Q'=PQ$, $K'=PK$ và $V'=PV$, nên $S'=PSP^\top$.

\subsection*{Hoán vị đi qua softmax theo hàng}
Chọn quy ước $(PX)_{i,:}=X_{\tau(i),:}$ với $\tau$ là một song ánh trên $\{1,\ldots,n\}$. Khi đó $(PSP^\top)_{ij}=S_{\tau(i),\tau(j)}$ và
\begin{align*}
  [\softmax_{\mathrm{row}}(PSP^\top)]_{ij}
    &=\frac{\exp(S_{\tau(i),\tau(j)})}{\sum_{a=1}^{n}\exp(S_{\tau(i),\tau(a)})}\\
    &=\frac{\exp(S_{\tau(i),\tau(j)})}{\sum_{b=1}^{n}\exp(S_{\tau(i),b})}\\
    &=[P\softmax_{\mathrm{row}}(S)P^\top]_{ij}.
\end{align*}
Đẳng thức thứ hai đúng vì $\tau(a)$ đi qua đúng một lần mỗi chỉ số cột. Do $P^\top P=I_n$,
\begin{align*}
  \Attn(PX)
    &=\softmax_{\mathrm{row}}(PSP^\top)PV\\
    &=P\softmax_{\mathrm{row}}(S)P^\top PV
      =P\softmax_{\mathrm{row}}(S)V=P\Attn(X).
\end{align*}
Trong multi-head attention, lập luận trên đúng cho từng đầu. Ghép đầu theo chiều đặc trưng rồi nhân bên phải với $W_O$ vẫn cho đúng phép hoán vị hàng: ghép các $PH_i$ cho kết quả $P$ nhân với ma trận ghép các $H_i$.

\subsection*{Giới hạn khi có mặt nạ hoặc mã hóa vị trí}
Với mặt nạ cộng $M\in(\R\cup\{-\infty\})^{n\times n}$ có ít nhất một vị trí không bị che ở mỗi hàng, đặt $\Attn_M(X)=\softmax_{\mathrm{row}}(S+M)V$. Phép hoán vị cột vẫn giữ tổng mẫu số (với $e^{-\infty}=0$), do đó
\[
  \Attn_{PMP^\top}(PX)=P\Attn_M(X),
\]
trong đó $PMP^\top$ được hiểu là đổi chỉ số hàng và cột của mặt nạ. Điều này không cho phép giữ nguyên một causal mask tùy ý. Chẳng hạn, lấy hai token vô hướng, $X=(0,2)^\top$, $W_Q=W_K=0$, $W_V=1$ và $P$ đổi chỗ hai hàng. Giữ mặt nạ nhân quả (hàng thứ nhất nhìn chính nó; hàng thứ hai nhìn cả hai) cho kết quả
\[
  \Attn_M(X)=(0,1)^\top,\qquad
  \Attn_M(PX)=(2,1)^\top\ne(1,0)^\top=P\Attn_M(X).
\]
Khi cộng một bảng mã hóa vị trí cố định $E$ vào biểu diễn token, đổi thứ tự token nhưng giữ vị trí tương ứng với $PX+E$, nói chung khác $P(X+E)$. Vì thế định lý về attention không PE không suy ra đẳng biến đối với thao tác giữ nguyên bảng vị trí này.

\section{Ma trận quay khối của mã hóa vị trí hình sin}
\label{sec:pe}
\textbf{Mệnh đề.} Xét vector cột $\PE(p)\in\R^{2m}$ gồm $m\ge1$ cặp đầy đủ theo thứ tự sin trước, cos sau. Với tần số cố định $\omega_i$, $i=0,\ldots,m-1$, đặt $\PE(p)_{2i}=\sin(p\omega_i)$ và $\PE(p)_{2i+1}=\cos(p\omega_i)$. Khi đó tồn tại $M_\Delta\in\R^{2m\times 2m}$ chỉ phụ thuộc độ dịch $\Delta$ (với hệ tần số đã cố định) sao cho $\PE(p+\Delta)=M_\Delta\PE(p)$ với mọi $p$. Trong Transformer gốc, $\omega_i=10000^{-2i/d_{\mathrm{model}}}$ \cite[§3.5]{attention}.

\subsection*{Chứng minh cho một cặp tọa độ}
Đặt $a=p\omega_i$ và $b=\Delta\omega_i$. Công thức cộng góc cho
\begin{align*}
  \sin(a+b)&=\cos b\,\sin a+\sin b\,\cos a,\\
  \cos(a+b)&=-\sin b\,\sin a+\cos b\,\cos a.
\end{align*}
Vì vậy, nếu $u_i(p)=(\sin(p\omega_i),\cos(p\omega_i))^\top$, ta có $u_i(p+\Delta)=M_{\Delta,i}u_i(p)$ với
\[
  M_{\Delta,i}=
    \begin{pmatrix}
      \cos(\Delta\omega_i)&\sin(\Delta\omega_i)\\
      -\sin(\Delta\omega_i)&\cos(\Delta\omega_i)
    \end{pmatrix}.
\]
Hai cột của khối này có chuẩn $1$ vì $\cos^2 b+\sin^2 b=1$, và tích vô hướng bằng $\cos b\sin b-\sin b\cos b=0$. Định thức bằng $\cos^2 b+\sin^2 b=1$. Do đó đây là khối quay; theo quy ước vector cột $(\sin,\cos)^\top$, nó là ma trận quay góc $-b$ trong mặt phẳng tọa độ chuẩn.

\subsection*{Ghép các cặp thành vector đầy đủ}
Lấy $M_\Delta=\diag(M_{\Delta,0},\ldots,M_{\Delta,m-1})$. Vì các khối ngoài đường chéo bằng $0$, phép nhân lên từng cặp tọa độ chính là $u_i(p)\mapsto u_i(p+\Delta)$. Ghép các cặp theo thứ tự ban đầu được $M_\Delta\PE(p)=\PE(p+\Delta)$. Các phần tử của $M_\Delta$ chỉ chứa $\Delta$ và những tần số đã cố định, không chứa $p$. Điều này cho thấy độ dịch cố định có biểu diễn tuyến tính trên PE; không khẳng định mô hình đã học được cách khai thác độ dịch.

\subsection*{Phạm vi: mệnh đề cần số chiều chẵn}
Giả thiết ``$m$ cặp đầy đủ'' là điều kiện thực sự cần, không phải quy ước trình bày. Cài đặt \texttt{SinusoidalPositionEncoding} trong \texttt{src/attention.py} của nhóm chấp nhận $d_{\mathrm{model}}$ lẻ: với $d_{\mathrm{model}}=2m+1$, mã sinh $m+1$ tần số $\omega_0,\ldots,\omega_m$, đặt $\sin$ vào $m+1$ cột chỉ số chẵn và $\cos$ vào $m$ cột chỉ số lẻ, nên
\[
  \PE(p)=\bigl(\sin p\omega_0,\ \cos p\omega_0,\ \ldots,\
                \sin p\omega_{m-1},\ \cos p\omega_{m-1},\
                \sin p\omega_m\bigr)^\top .
\]
Tọa độ cuối $\sin(p\omega_m)$ không có $\cos(p\omega_m)$ đi kèm. Trong trường hợp này không tồn tại $M_\Delta$ nào thỏa mệnh đề, trừ khi $\sin(\Delta\omega_m)=0$.

\textbf{Bổ đề (độc lập tuyến tính).} Cho các tần số phân biệt $\omega_0,\ldots,\omega_m\in(0,\pi)$. Nếu $L\ge 2m+2$ thì $2m+2$ hàm $p\mapsto\sin(p\omega_i)$ và $p\mapsto\cos(p\omega_i)$, $0\le i\le m$, độc lập tuyến tính trên $\R$ khi xét như hàm trên $\{0,1,\ldots,L-1\}$.

\emph{Chứng minh.} Đặt $z_i=e^{\mathrm i\omega_i}$. Mỗi tổ hợp tuyến tính thực của các hàm trên viết được thành $\sum_{i=0}^{m}\bigl(\gamma_i z_i^{\,p}+\overline{\gamma_i}\,\overline{z_i}^{\,p}\bigr)$. Vì $\omega_i\in(0,\pi)$ phân biệt, $2m+2$ số $z_0,\overline{z_0},\ldots,z_m,\overline{z_m}$ đôi một khác nhau. Nếu tổng này triệt tiêu tại $p=0,\ldots,2m+1$ thì hệ Vandermonde với các nút phân biệt đó khả nghịch, buộc mọi $\gamma_i=0$.\hfill$\square$

Với lựa chọn của Transformer gốc, $\omega_i=10000^{-2i/d_{\mathrm{model}}}\in(0,1]\subset(0,\pi)$ và đôi một phân biệt, nên bổ đề áp dụng được khi $\texttt{max\_len}\ge d_{\mathrm{model}}+1$.

\textbf{Suy ra điều phải bác bỏ.} Giả sử có $M_\Delta\in\R^{(2m+1)\times(2m+1)}$ với $\PE(p+\Delta)=M_\Delta\PE(p)$ cho mọi $p$ trong miền vị trí. Đọc hàng cuối (chỉ số $2m$) và dùng công thức cộng góc:
\[
  \cos(\Delta\omega_m)\sin(p\omega_m)+\sin(\Delta\omega_m)\cos(p\omega_m)
  =\sum_{i=0}^{m-1}\bigl(a_i\sin(p\omega_i)+b_i\cos(p\omega_i)\bigr)
   +a_m\sin(p\omega_m),
\]
trong đó $a_i,b_i$ là các phần tử của hàng đó, không phụ thuộc $p$. Chuyển vế:
\[
  \sin(\Delta\omega_m)\cos(p\omega_m)
  =\sum_{i=0}^{m-1}\bigl(a_i\sin(p\omega_i)+b_i\cos(p\omega_i)\bigr)
   +\bigl(a_m-\cos(\Delta\omega_m)\bigr)\sin(p\omega_m).
\]
Theo bổ đề, hàm $\cos(p\omega_m)$ độc lập tuyến tính với tất cả các hàm ở vế phải, nên đẳng thức chỉ có thể đúng với mọi $p$ khi $\sin(\Delta\omega_m)=0$, tức $\Delta\omega_m\in\pi\mathbb Z$. Vì $0<\omega_m\le 1$, điều này loại mọi độ dịch nguyên $\Delta$ với $|\Delta|<\pi/\omega_m$, trong đó có $\Delta=1$. Vậy mệnh đề thất bại ở số chiều lẻ.

\textbf{Hệ quả cho cài đặt.} Thí nghiệm A.2 dùng $d_{\mathrm{model}}=64$ (chẵn), nên bảng PE được dùng thực tế vẫn thỏa mệnh đề; phần kiểm thử $d_{\mathrm{model}}$ lẻ của nhóm chỉ đối chiếu giá trị bảng, không kiểm tính chất dịch vị trí. Muốn giữ tính chất cho số chiều lẻ, cần bỏ tọa độ $\sin$ lẻ cặp hoặc đệm lên số chiều chẵn; đây là lựa chọn thiết kế, không phải hệ quả của chứng minh trên.

\section{Khởi tạo, hạng, số tham số và hợp nhất LoRA}
\label{sec:lora}
\textbf{Thiết lập.} Xét lớp tuyến tính có $W\in\R^{d\times k}$, với $d,k\ge1$, và đầu vào cột $h\in\R^k$. LoRA đóng băng $W$, học $A\in\R^{r\times k}$ và $B\in\R^{d\times r}$, với $r\ge1$ nguyên và $\alpha$ cố định. Không xét huấn luyện bias; nếu có bias cố định $b\in\R^d$, nó được cộng như nhau trong cả hai cách tính. Đặt $\Delta W=(\alpha/r)BA$ và $W'=W+\Delta W$ \cite[§4.1–4.2]{lora}.

\subsection*{Khởi tạo không thay đổi lớp gốc}
Khi $B=0_{d\times r}$, mỗi phần tử $(BA)_{ij}=\sum_{a=1}^{r} B_{ia}A_{aj}$ bằng $0$, bất kể giá trị khởi tạo của $A$. Do đó $\Delta W=0_{d\times k}$ và $W'=W$. Bài báo LoRA chọn $A$ khởi tạo ngẫu nhiên và $B=0$; kết quả $W'=W$ chỉ cần điều kiện thứ hai.

\subsection*{Hạng của cập nhật bị chặn bởi $r$}
Xem $A$ là ánh xạ từ $\R^k$ vào $\R^r$ và $B$ là ánh xạ từ $\R^r$ vào $\R^d$. Với mọi $h$, $BAh=B(Ah)\in\im(B)$. Vì $\im(B)$ sinh bởi $r$ cột của $B$,
\[
  \rank(BA)=\dim\im(BA)\le\dim\im(B)=\rank(B)\le r.
\]
Suy ra $\rank(\Delta W)\le r$. Mệnh đề này chặn hạng của \emph{phần cập nhật}, không chặn hạng của $W'$.

\subsection*{Đếm tham số thực sự được huấn luyện}
Học trực tiếp toàn bộ $W$ cần cập nhật $dk$ hệ số. Với LoRA, $A$ có $rk$ phần tử và $B$ có $dr$ phần tử. Vì $W$, $\alpha$ và bias không được huấn luyện, số tham số cập nhật là $rk+dr=r(d+k)$. Số này nhỏ hơn $dk$ khi $r<dk/(d+k)$; tỷ lệ chính xác so với học cả $W$ là $r(d+k)/(dk)$. Đây không phải tổng dung lượng mô hình, vì vẫn cần trọng số gốc.

\subsection*{Hợp nhất không thêm phép tính trong lan truyền xuôi}
Trong suy luận xác định với adapter đã cố định, hai nhánh cho cùng ánh xạ như một lớp tuyến tính duy nhất:
\[
  Wh+\frac{\alpha}{r}B(Ah)=\left(W+\frac{\alpha}{r}BA\right)h=W'h.
\]
Tính và lưu $W'$ một lần trước khi phục vụ. Sau đó mỗi đầu vào dùng đúng một phép nhân ma trận--vector kích thước $d\times k$ như ở lớp gốc: $dk$ phép nhân và $d(k-1)$ phép cộng nếu tính trực tiếp, cộng cùng chi phí bias nếu có. Không còn các phép nhân phụ $Ah$, $B(Ah)$ hay phép cộng hai nhánh trên mỗi đầu vào. Kết luận nói về chi phí số học \emph{sau hợp nhất}, giả sử cùng kiểu thực thi lớp tuyến tính. Nó không bao gồm chi phí tạo $W'$, không bảo đảm khớp bit trong số học dấu phẩy động, và không tự mở rộng thành kết luận về độ trễ khi đổi kernel hay tái lượng tử hóa.

\section{Nghiệm tối ưu có phạt KL và hàm mất mát DPO}
\label{sec:dpo}
\textbf{Thiết lập.} Cố định câu hỏi $x$. Ký hiệu $R(x,y)$ là reward của câu trả lời $y$ (đề dùng $r(x,y)$), $\beta>0$, và $\pref(\cdot\mid x)$ là phân phối tham chiếu cố định. Trước hết xét tập câu trả lời hữu hạn, không rỗng $\mathcal Y_x$, với reward hữu hạn trên hỗ trợ $S_x=\{y\in\mathcal Y_x:\pref(y\mid x)>0\}$. Tối ưu trên mọi phân phối $\pi(\cdot\mid x)$ có hỗ trợ nằm trong $S_x$. Chính sách gán khối lượng dương cho $y\notin S_x$ có KL vô hạn và không thể là nghiệm hữu hạn; nếu xét trên toàn bộ $\mathcal Y_x$, giả sử reward hữu hạn trên toàn bộ tập này. Quy ước hạng tử có hệ số xác suất bằng $0$ trong KL bằng $0$. Bài toán và bước đưa vào Bradley–Terry theo \cite[§3–4, phương trình (1)–(7), Phụ lục A.1–A.2]{dpo}.

\subsection*{Xây dựng một phân phối ứng viên}
Viết tắt $\pi_y=\pi(y\mid x)$, $q_y=\pref(y\mid x)$ và $R_y=R(x,y)$ trên $S_x$. Hàm mục tiêu là
\[
  J_x(\pi)=\sum_{y\in S_x}\pi_y R_y
    -\beta\sum_{y\in S_x:\,\pi_y>0}\pi_y\log\frac{\pi_y}{q_y}.
\]
Đặt $Z(x)=\sum_{y\in S_x}q_y\exp(R_y/\beta)$. Do tập hữu hạn, hỗ trợ không rỗng và các số hạng dương hữu hạn, $0<Z(x)<\infty$. Định nghĩa
\begin{equation}
  \pi^*(y\mid x)=\frac{\pref(y\mid x)\exp(R(x,y)/\beta)}{Z(x)}
  \quad (y\in S_x),
  \label{eq:gibbs}
\end{equation}
và đặt $\pi^*(y\mid x)=0$ ngoài $S_x$. Tổng tử số trên $S_x$ chính bằng $Z(x)$, nên $\pi^*$ là một phân phối, dương trên toàn bộ $S_x$.

\subsection*{Chứng minh tối ưu toàn cục và duy nhất}
Ký hiệu $\pi_y^*=\pi^*(y\mid x)$. Từ \eqref{eq:gibbs}, $\log(\pi_y^*/q_y)=R_y/\beta-\log Z(x)$. Với quy ước bỏ các hạng có $\pi_y=0$ trong tổng chứa log, ta được
\begin{align*}
  \beta\KL(\pi\|\pi^*)
     &=\beta\sum_{y:\,\pi_y>0}\pi_y\log\frac{\pi_y}{\pi_y^*}\\
     &=\beta\sum_{y:\,\pi_y>0}\pi_y\log\frac{\pi_y}{q_y}
       -\sum_y\pi_y R_y+\beta\log Z(x)\\
     &=\beta\log Z(x)-J_x(\pi).
\end{align*}
Để thấy điều kiện dấu bằng, dùng $-\log t\ge 1-t$ với $t>0$, dấu bằng chỉ tại $t=1$. Với $T=\{y:\pi_y>0\}$, áp dụng cho $t=\pi_y^*/\pi_y$:
\[
  \KL(\pi\|\pi^*)\ge\sum_{y\in T}(\pi_y-\pi_y^*)
       =1-\sum_{y\in T}\pi_y^*\ge0.
\]
Nếu KL bằng $0$, bất đẳng thức đầu buộc $\pi_y^*=\pi_y$ trên $T$, bất đẳng thức sau buộc $\pi^*$ không gán khối lượng ngoài $T$. Ngược lại, $\pi=\pi^*$ cho KL bằng $0$. Vì $\beta>0$, $J_x(\pi)$ đạt cực đại $\beta\log Z(x)$ đúng tại $\pi^*$; không còn nghiệm nào khác, kể cả ở biên của miền phân phối.

Với miền câu trả lời đếm được, một điều kiện đủ để dùng nguyên lập luận là reward bị chặn trên hỗ trợ theo cả hai phía: nếu $a\le R_y\le b$ thì $0<e^{a/\beta}\le Z(x)\le e^{b/\beta}<\infty$. Khi ấy kỳ vọng reward hữu hạn và $\log(\pi_y^*/q_y)$ bị chặn. Đẳng thức vẫn đúng với KL nhận giá trị $+\infty$, cho mục tiêu $-\infty$ chứ không tạo hiệu $\infty-\infty$. Nếu reward không bị chặn, cần kiểm tra riêng việc chuẩn hóa và tính xác định của mục tiêu; không suy ra các điều kiện này chỉ từ việc $\pref$ là một phân phối.

\subsection*{Từ Bradley–Terry tới xác suất ưu tiên}
Xét cặp $y_w,y_l\in S_x$ cho cùng $x$, trong đó $y_w$ là đáp án được ưu tiên và $y_l$ là đáp án còn lại. Mô hình Bradley–Terry gán xác suất
\begin{align*}
  \Pr(y_w\succ y_l\mid x,y_w,y_l)
   &=\frac{\exp(R(x,y_w))}{\exp(R(x,y_w))+\exp(R(x,y_l))}\\
   &=\sigma\bigl(R(x,y_w)-R(x,y_l)\bigr),
   \qquad \sigma(z)=\frac{1}{1+e^{-z}}.
\end{align*}
Giải \eqref{eq:gibbs} theo reward:
\[
  R(x,y)=\beta\log\frac{\pi^*(y\mid x)}{\pref(y\mid x)}+\beta\log Z(x).
\]
Lấy hiệu giữa hai đáp án cùng $x$, hai hạng $\beta\log Z(x)$ triệt tiêu. Do đó xác suất ưu tiên chỉ còn phụ thuộc vào hiệu hai log-tỷ số chính sách tối ưu/tham chiếu. Bước triệt tiêu không áp dụng cho hai đáp án của hai câu hỏi khác nhau.

\subsection*{Hợp lý cực đại theo chính sách tham số hóa}
Dùng $\pi_\theta$ thay cho chính sách chưa biết $\pi^*$ trong xác suất ưu tiên. Hai đáp án của mỗi cặp phải có xác suất dương dưới cả $\pi_\theta$ và $\pref$ để log-tỷ số hữu hạn. Đặt
\begin{equation}
  z_\theta(x,y_w,y_l)
   =\beta\left[
      \log\frac{\pi_\theta(y_w\mid x)}{\pref(y_w\mid x)}
      -\log\frac{\pi_\theta(y_l\mid x)}{\pref(y_l\mid x)}
   \right].
  \label{eq:dpo-logit}
\end{equation}
Với tập $\mathcal D$ các bộ ba đã sắp đáp án theo nhãn ưu tiên, mỗi quan sát đóng góp xác suất $\sigma(z_\theta)$ vào hợp lý của nhãn, có điều kiện trên câu hỏi và cặp đáp án được đưa ra. Với các quan sát độc lập có điều kiện theo mô hình nhãn, hợp lý là tích các xác suất đó. Lấy log biến tích thành tổng; đổi dấu và lấy trung bình trên $\mathcal D$ cho
\begin{equation}
  \boxed{\mathcal L_{\mathrm{DPO}}(\theta)
    =-\E_{(x,y_w,y_l)\sim\mathcal D}\bigl[\log\sigma(z_\theta(x,y_w,y_l))\bigr].}
  \label{eq:dpo-loss}
\end{equation}
Đây là mất mát DPO, không cần tính $Z(x)$ hay huấn luyện một reward model riêng để tính mất mát. Có thể kiểm dấu: nếu hai log-tỷ số trong \eqref{eq:dpo-logit} bằng nhau, mất mát một cặp bằng $-\log(1/2)=\log 2$; nếu hiệu tăng, $\sigma$ tăng và mất mát giảm. Tính duy nhất ở trên là tính duy nhất của $\pi^*$ với reward đã cho khi tối ưu trên các phân phối, không phải tính duy nhất của tham số $\theta$ khi học từ cặp hữu hạn.

\section{Cross-entropy kỳ vọng, độ phân kỳ KL và perplexity}
\label{sec:ppl}
\textbf{Thiết lập.} Xét phân phối dữ liệu cố định $\pdata$ và một họ phân phối mô hình $p_\theta$, trên cùng một miền rời rạc hữu hạn $\mathcal Y$. Giả sử $p_\theta(y)>0$ với mọi $y$ sao cho $\pdata(y)>0$ khi cần giá trị hữu hạn. Các hạng có $\pdata(y)=0$ được bỏ khỏi tổng; nếu $\pdata(y)>0$ nhưng $p_\theta(y)=0$, cả cross-entropy và KL đều là $+\infty$.

\subsection*{Phân rã cross-entropy kỳ vọng}
Đặt $S=\{y:\pdata(y)>0\}$. Theo định nghĩa, cross-entropy của dữ liệu dưới mô hình là kỳ vọng của log-xác suất âm, còn entropy dữ liệu đo cùng đại lượng dưới chính phân phối dữ liệu:
\begin{align*}
  H(\pdata,p_\theta)&=-\sum_{y\in S}\pdata(y)\log p_\theta(y),\\
  H(\pdata)&=-\sum_{y\in S}\pdata(y)\log\pdata(y).
\end{align*}
Khai triển tỷ số trong KL thành hiệu hai logarit:
\begin{align*}
  \KL(\pdata\|p_\theta)
   &=\sum_{y\in S}\pdata(y)\log\frac{\pdata(y)}{p_\theta(y)}\\
   &=-H(\pdata)+H(\pdata,p_\theta).
\end{align*}
Vì $H(\pdata)$ hữu hạn và không phụ thuộc $\theta$, hai mục tiêu khác nhau đúng một hằng số: chúng xếp thứ tự các mô hình như nhau và có cùng tập nghiệm cực tiểu (nếu đạt cực tiểu):
\[
  \argmin_\theta H(\pdata,p_\theta)
    =\argmin_\theta \KL(\pdata\|p_\theta).
\]
Trên miền đếm được, lập luận trừ hằng số vẫn dùng được khi $H(\pdata)$ hữu hạn và cross-entropy/KL được hiểu trong miền giá trị mở rộng. Không được lấy hiệu $\infty-\infty$ nếu entropy dữ liệu vô hạn.

\subsection*{Phiên bản có ngữ cảnh cho dự đoán token}
Ký hiệu $C$ là ngữ cảnh và $Y$ là token đích. Giữ phân phối dữ liệu của $C$ cố định, áp dụng đẳng thức trên cho mỗi ngữ cảnh $c$ và lấy kỳ vọng:
\begin{align*}
  \E_{(C,Y)\sim\pdata}[-\log p_\theta(Y\mid C)]
   &= H_{\pdata}(Y\mid C)\\
   &\quad+\E_{C\sim\pdata}\bigl[\KL(\pdata(\cdot\mid C)\|p_\theta(\cdot\mid C))\bigr],
\end{align*}
trong đó $H_{\pdata}(Y\mid C)=\E_{C\sim\pdata}[H(\pdata(\cdot\mid C))]$ không phụ thuộc mô hình. Vì từ vựng hữu hạn, entropy có điều kiện hữu hạn. Đây là phát biểu tương ứng khi mô hình ngôn ngữ được huấn luyện dự đoán token tiếp theo theo ngữ cảnh đúng từ dữ liệu.

\subsection*{Từ log-xác suất trung bình theo token đến PPL}
Trên tập đánh giá, gọi $(c_t,y_t)$ là ngữ cảnh và token đích ở vị trí được chấm thứ $t$, với tổng cộng $N>0$ token được chấm. Với logarit tự nhiên và xác suất dương tại các vị trí này,
\begin{equation}
  \bar H=-\frac{1}{N}\sum_{t=1}^{N}\log p_\theta(y_t\mid c_t),\qquad
  \PPL=e^{\bar H}
    =\left(\prod_{t=1}^{N} p_\theta(y_t\mid c_t)\right)^{-1/N}.
  \label{eq:ppl}
\end{equation}
Đẳng thức cuối dùng $\exp(\sum_t\log a_t)=\prod_t a_t$. Như vậy PPL là nghịch đảo trung bình nhân của xác suất token đúng; vì hàm mũ tăng chặt, cross-entropy theo token thấp hơn tương đương PPL thấp hơn trên cùng thiết lập chấm. Ví dụ minh họa: hai token có xác suất đúng lần lượt $1/2$ và $1/8$ cho $\bar H=\log 4$ và $\PPL=4$.

Khi cộng qua nhiều batch hoặc nhiều câu, cộng tổng log-xác suất âm rồi chia cho tổng số token được chấm, không lấy trung bình các PPL riêng lẻ. Loại các vị trí padding khỏi cả tử và mẫu; khi so sánh, dùng cùng tokenizer, dữ liệu và quy ước chấm token. Ngược lại, $\log\pi_\theta(y\mid x)$ trong \eqref{eq:dpo-logit} là log-xác suất \emph{cả câu trả lời}. Theo quy tắc dây chuyền nó bằng $\sum_{t=1}^{T}\log\pi_\theta(y_t\mid x,y_{<t})$, với $T$ là số token của đáp án theo cùng quy ước kết thúc chuỗi cho hai mô hình. Không tính log-xác suất của prompt cố định và không chia tổng này cho $T$ trong DPO chuẩn.

\begin{thebibliography}{9}
\bibitem{de}
Phạm Văn Khánh. \emph{Danh mục 20 bài tập lớn: Học máy cơ bản, mô hình ngôn ngữ lớn, học tăng cường và tích hợp}. Bản ngày 12/11/2026, 46 trang; Đề tài 9, Phần B, tr.~18.
\bibitem{attention}
Vaswani, A. và cộng sự. \emph{Attention Is All You Need}. NeurIPS 2017; bản \href{https://arxiv.org/abs/1706.03762v7}{arXiv:1706.03762v7}.
\bibitem{lora}
Hu, E. J. và cộng sự. \emph{LoRA: Low-Rank Adaptation of Large Language Models}. ICLR 2022; bản \href{https://arxiv.org/abs/2106.09685v2}{arXiv:2106.09685v2}.
\bibitem{dpo}
Rafailov, R. và cộng sự. \emph{Direct Preference Optimization: Your Language Model Is Secretly a Reward Model}. NeurIPS 2023; bản \href{https://arxiv.org/abs/2305.18290v3}{arXiv:2305.18290v3}.
\end{thebibliography}
\end{document}
